% Cote 119: offprints, reports, courses, CV and letters (typed and printed part): archive cote 119, site page 171 (PDF page 172).
% typescript. Draft compared: none.
% AI-transcribed by Claude Opus 5.5 (claude-opus-5-5, adaptive thinking, effort high), Message Batches API, batch r119_full (msgbatch_01DwZNm5Y7YsKBWyAN4sFs8w), 2026-11-13; not yet checked by eye.
\scanpage{119/171}
\mspage{- 10 -}

$R_1$, $R_2$, $R_3$, $R_4$. On vérifie (comme en 3.4) que
\[ \alpha(R) = \int_\gamma f(Z)\, dZ = \sum_{i=1}^{4} \int_{\gamma_i} f(Z)\, dZ = \sum_{i=1}^{4} \alpha(R_i) \]
Donc parmi ces quatre rectangles il y en a au moins un tel que
$|\alpha(R_i)| \geq \frac{1}{4}\, |\alpha(R)|$ . Appelons $R^{(1)}$ ce \unsure{rectangme|rectangle}.
En réitérant l'opération précédente sur le rectangle $R^{(1)}$ on trouve un
rectangle $R^{(2)}$ tel que $|\alpha(R^{(2)})| \geq \frac{1}{4}\, |\alpha(R^{(1)}| \geq \frac{1}{4^2}\, |\alpha(R)|$.
On peut recommencer indéfiniment la même opération. On obtiendra ainsi une suite
de rectangles emboités ; le k-ième aura ses côtés $2^k$ fois plus petit que ceux
de $R$ et son aire sere $4^k$ fois plus petite que celle de $R$.
Si $\gamma(k)$ désigne le bord orienté du rectangle $R^{(k)}$ on a :
\[ \text{(4.3.1.)} \qquad \Big|\int_{\gamma(k)} f(Z)\, dZ\Big| \geq \frac{1}{4^k}\, \alpha(R) \ . \]
L'intersection des rectangles $R^{(k)}$ est égale à un point $Z_0$
(justifiez). On a $Z_0 \in U$ . Comme $f$ est holomorphe en $Z_0$. On a :
\[ f(Z) = f(Z_0) + f'(Z_0)\ (Z-Z_0) + \epsilon(Z)\, |Z-Z_0| \quad \text{avec} \ \lim_{Z \to Z_0} \epsilon(Z) = 0 \]
on en désuit :
\[ \int_{\gamma(k)} f(Z)\, dZ = f(Z_o) \int_{\gamma(k)} dZ + f'(Z_0) \int_{\gamma(k)} (Z-Z_0)\, \unsure{dZ} + \int_{\gamma(k)} \epsilon(Z)\, |Z-Z_0|\, dZ \]
les deux premières intégrales sont nulles (proposition 3.3). Si $a$ et $b$ désignent
les longueurs des côtés du rectangle $R$ avec $a \geq b$. On a :
\[ |Z-Z_0| \leq \sqrt{a^2+b^2}\ .\ 2^{-k} \]
périmètre de $R^{(k)} = 2(a+b).2^{-k}$ d'iù si $k$ est assez grand pour que $|\epsilon(Z)| \leq \epsilon$
quelque soit $Z$ sur la frontière de $R^{(k)}$ :
(4.3.2)
\[ \Big|\int_{\gamma(k)} \epsilon(Z)\, |Z-Z_0|\, dZ\Big| \leq \epsilon\ .\ K\ .\ 4^{-k} \quad \text{où $K$ est une constante positive.} \]
On a donc en comparant (4.3.1) et (4.3.2) :
$\alpha(R) \leq \epsilon.K$ et ce, pour tout,$\epsilon > 0$, ce qui montre que $\alpha(R) = 0$ et
achève la démonstration.

\margin{171}
